\section{Step 3: Construct the local resurgent basis}
\label{sec:exact-wkb}
% BEGIN CURATED SUBJECT INDEX
\index{Riccati equation}
% END CURATED SUBJECT INDEX

\calculationroute
  {The spectral curve $y^2=Q(x,E)$ and the asymptotic coefficients fixed in
  Sec.~\ref{sec:textbook-scattering}.}
  {Generate the all-order Riccati series, separate odd and even parts, Borel
  resum in a declared direction, and identify the turning-point Stokes data.}
  {Produce normalized sectorial WKB branches whose ambiguity and domain of
  validity are explicit.}

\subsection{Riccati equation and formal WKB solutions}
\label{subsec:wkb-riccati}
% BEGIN CURATED SUBJECT INDEX
\index{Riccati equation}
\index{exact WKB}
\index{monodromy}
\index{Langer correction}
% END CURATED SUBJECT INDEX

Consider the Schr\"odinger equation in the form
\begin{equation}
  \left[-\hbar^2\frac{\rmd^2}{\rmd x^2}+Q(x,E)\right]\psi(x)=0,
  \qquad
  Q(x,E)=2m\left[V(x)-E\right] .
  \label{eq:wkb-schrodinger}
\end{equation}
The coordinate $x$ and energy $E$ may be complex.  Introduce the logarithmic
derivative
\begin{equation}
  \psi(x)=\rme^{\int^x S(x',\hbar)\rmd x'} .
  \label{eq:wkb-log-ansatz}
\end{equation}
Then $S$ satisfies the Riccati equation
\begin{equation}
  S^2+\frac{\partial S}{\partial x}
  =\frac{Q}{\hbar^2} .
  \label{eq:wkb-riccati}
\end{equation}
Seek a formal solution
\begin{equation}
  S^{(\pm)}(x,\hbar)
  =\sum_{n=-1}^{\infty}\hbar^nS_n^{(\pm)}(x) .
  \label{eq:wkb-S-series}
\end{equation}
The leading terms are
\begin{align}
  S_{-1}^{(\pm)}&=\pm\sqrt{Q},
  &S_0^{(\pm)}&=-\frac{1}{4}\frac{\partial_xQ}{Q} .
  \label{eq:wkb-leading}
\end{align}
For $n\ge0$, the remaining coefficients obey
\begin{equation}
  2S_{-1}^{(\pm)}S_{n+1}^{(\pm)}
  =-\partial_xS_n^{(\pm)}
  -\sum_{j=0}^{n}S_j^{(\pm)}S_{n-j}^{(\pm)} .
  \label{eq:wkb-recursion}
\end{equation}

Define odd and even parts under exchange of the two leading branches:
\begin{align}
  \vsub{S}{odd}
  &=\frac{1}{2}\left(S^{(+)}-S^{(-)}\right),
  \notag\\
  \vsub{S}{even}
  &=\frac{1}{2}\left(S^{(+)}+S^{(-)}\right)
  =-\frac{1}{2}\partial_x\log \vsub{S}{odd} .
  \label{eq:wkb-odd-even}
\end{align}
The formal WKB solutions normalized at a reference point $x_0$ are
\begin{equation}
  \psi_{\pm}(x)
  =\frac{1}{\sqrt{\vsub{S}{odd}(x)}}
  \rme^{
  \pm\int_{x_0}^{x}\vsub{S}{odd}(x')\rmd x'
  } .
  \label{eq:formal-wkb-solutions}
\end{equation}
The square root, integration path, and sheet of $\sqrt Q$ are part of the
definition.  A local asymptotic series becomes a global spectral statement
only after these data and the Stokes transitions are specified.

If the original differential equation contains a first-derivative term, a
local gauge transformation first removes it and shifts $Q$ by a term of order
$\hbar^2$.  Such terms are easily dismissed in leading semiclassics but can
carry essential monodromy at a regular singular point.  The radial Langer term
in Section~\ref{sec:radial-wkb} is the principal example.

\subsection{Borel transform and lateral resummation}
\label{subsec:borel-resummation}
% BEGIN CURATED SUBJECT INDEX
\index{Borel transform}
\index{Borel summation}
\index{lateral Borel sum}
\index{Borel plane}
\index{Stokes automorphism}
\index{exact WKB}
\index{monodromy}
% END CURATED SUBJECT INDEX

The series in Eq.~\eqref{eq:wkb-S-series} is generally divergent.  For a
formal series
\begin{equation}
  f(\hbar)=\sum_{n=0}^{\infty}a_n\hbar^n,
  \label{eq:formal-series}
\end{equation}
define its Borel transform by
\begin{equation}
  \mathcal Bf(\xi)
  =\sum_{n=0}^{\infty}\frac{a_n}{n!}\xi^n .
  \label{eq:borel-transform}
\end{equation}
If $\mathcal Bf$ can be analytically continued with suitable growth along the
ray $\arg\xi=\varphi$, its directional Borel sum is
\begin{equation}
  \mathcal S_{\varphi}f(\hbar)
  =\frac{1}{\hbar}
  \int_{0}^{\rme^{\rmi\varphi}\infty}
  \rmd\xi\,
  \rme^{-\xi/\hbar}\mathcal Bf(\xi) .
  \label{eq:borel-sum}
\end{equation}
When a Borel singularity lies on the integration ray, define lateral sums
$\mathcal S_{\varphi\pm}$.  Their difference is exponentially small in
$\hbar$ and is governed by a Stokes automorphism.  This is the point at which
nonperturbative saddles enter a perturbative expansion.  In a metastable
problem, the same analytic discontinuity that resolves a Borel ambiguity
produces the imaginary part of the resonance energy
\cite{Bogomolny:1980ur,Zinn-Justin:1981qzi,Dunne:2002rq,
Dunne:2015eaa}.

The phrase \textit{exact WKB} does not mean truncating the WKB series at high
order.  It means using Borel-resummed WKB solutions and their exact connection
formulae.  Truncation may approximate a Borel sum before the least term, but it
does not determine lateral continuation or global monodromy.

\subsection{Turning points and Stokes geometry}
\label{subsec:stokes-geometry}
% BEGIN CURATED SUBJECT INDEX
\index{Stokes curve}
\index{turning point}
\index{monodromy}
% END CURATED SUBJECT INDEX

A simple turning point $a$ satisfies
\begin{equation}
  Q(a,E)=0,
  \qquad
  \partial_xQ(a,E)\ne0 .
  \label{eq:simple-turning-point}
\end{equation}
Choose a branch of the classical action
\begin{equation}
  W_a(x;E)=\int_a^x\sqrt{Q(x',E)}\rmd x' .
  \label{eq:classical-action}
\end{equation}
In the convention used here, Stokes curves are trajectories on which
\begin{equation}
  \im\left[\frac{W_a(x;E)}{\hbar}\right]=0 .
  \label{eq:stokes-curve}
\end{equation}
Some literature interchanges the names Stokes and anti-Stokes curves; the
connection matrices, rather than the label, fix the convention.  Three curves
emanate from a simple turning point.  On crossing one of them, the coefficient
of a subdominant solution changes by a multiple of the dominant solution.

Let the local solution be represented by a coefficient column
$\bm c=(c_+,c_-)^{\mathsf T}$.  Up to orientation and normalization, simple
turning-point connections have triangular form
\begin{align}
  M_+(\sigma)&=
  \begin{pmatrix}1&\sigma\\0&1\end{pmatrix},
  &
  M_-(\sigma)&=
  \begin{pmatrix}1&0\\\sigma&1\end{pmatrix},
  \label{eq:stokes-matrices}
\end{align}
where the Stokes multiplier $\sigma$ is $+\rmi$ or $-\rmi$ for the standard
Airy normalization, depending on the crossing direction.  A global connection
problem is an ordered product of such matrices, propagation factors, and
monodromy matrices.  The physical boundary condition states that one component
of the final coefficient column vanishes.

\subsection{Voros symbols and exact quantization}
\label{subsec:voros-symbols}
% BEGIN CURATED SUBJECT INDEX
\index{Jost function}
\index{connection coefficient}
\index{boundary condition!incoming and outgoing}
\index{Borel summation}
\index{Voros symbol}
% END CURATED SUBJECT INDEX

For a path or cycle $\gamma$ on the spectral curve $y^2=Q(x,E)$, define the
formal Voros integral and symbol by
\begin{align}
  \VorosPeriod_{\gamma}(E,\hbar)
  &=\int_{\gamma}\vsub{S}{odd}(x,\hbar)\rmd x,
  \notag\\
  \VorosExp_{\gamma}(E,\hbar)
  &=\rme^{\VorosPeriod_{\gamma}(E,\hbar)} .
  \label{eq:voros-symbol}
\end{align}
Both are understood through a specified lateral Borel sum.  For a simple
two-turning-point bound-state problem, the connection condition often reduces
to
\begin{equation}
  1+\VorosExp_{\gamma}(E,\hbar)=0 .
  \label{eq:simple-exact-quantization}
\end{equation}
At leading order this gives the Bohr--Sommerfeld rule including the Maslov
phase.  More complicated graphs produce polynomial or cluster-like relations
among several Voros symbols
\cite{Voros:1983xx,Delabaere:1999xx,Iwaki:2014vad}.

For a resonance, the algebra is the same but the selected subdominant sectors
are different.  Let $\Cincoming(E,\hbar)$ be the connection coefficient
of the incoming asymptotic solution after transporting an outgoing solution
across the Stokes graph.  The exact resonance condition is
\begin{equation}
  \Cincoming(E,\hbar)=0 .
  \label{eq:exact-resonance-condition}
\end{equation}
Thus the last arrow in Eq.~\eqref{eq:intro-dictionary} is not an analogy:
$\Cincoming$ is a Jost function expressed in resurgent coordinates.

\subsection{Local normalization versus global invariants}
\label{subsec:wkb-gauge}
% BEGIN CURATED SUBJECT INDEX
\index{boundary condition!incoming and outgoing}
\index{Voros symbol}
\index{monodromy}
% END CURATED SUBJECT INDEX

The starting point of a WKB integral, the normalization of local bases, and
the detailed shape of a connection path are not separately physical.  Changing
them conjugates connection matrices or multiplies rows and columns by nonzero
factors.  Zeros of the incoming coefficient, closed Voros periods, and
monodromy conjugacy classes are invariant.  We will refer to the local choices
as a connection gauge.  This language is useful in radial problems, where an
open path from the origin to infinity and a closed contour encircling the
origin can encode identical quantization data
\cite{Morikawa:2025ezl}.
